Este documento reúne, en un solo texto, lo que el equipo ha escrito hasta ahora
sobre el proyecto \textit{Escanea y Entra: Prototipo de Acceso QR para el
TESCHA}, visto desde la gestión de proyectos. Unifica la ficha inicial del
proyecto, el plan de calidad, el tablero de metas y seguimiento, la
presentación del avance y el documento de investigación del equipo.

El proyecto busca probar, en un solo acceso del plantel, si un código QR en el
celular del estudiante vuelve más rápida y más confiable la entrada. Las
secciones siguientes explican por qué se hace, qué se va a construir, hasta
dónde llega, quién hace cada parte, cuánto tiempo y dinero se prevé, cómo se
cuidará la calidad y cómo se dará seguimiento. Las partes técnicas se explican
con palabras sencillas.

Algunos datos todavía no existen porque el proyecto apenas arranca. Cuando
una cifra depende de las mediciones del diagnóstico o del piloto, el documento
la deja como pendiente y no la estima.
